\section{妙鸭：把 AI 用得好的互联网产品}

本节从个人经历转入产品案例。理解妙鸭，是理解整期访谈的第一道门：它既是张月光最广为人知的成功案例，也是他后来反思“什么才是 AI Native”的反例。

妙鸭的成功来自几个产品判断：图像技术当时相对成熟；不要做 Midjourney 式通用图像生成；选足够垂直的写真人像；避开电商长链条，选择 2C 可定价的写真生意；把效果目标压成“真、像、美”。这些判断都很优秀，但张月光后来认为，妙鸭仍然是互联网产品思维。

\begin{figure}[H]
\centering
\includegraphics[width=0.94\textwidth]{miaoya-not-ai-native.png}
\caption{妙鸭为什么不是 AI Native：固定输入、固定链路、固定输出。}
\end{figure}

\begin{knowledgebox}{读图：妙鸭的稳定来自限制自由度}
图中三段分别是输入固定、链路固定、输出固定。用户上传约 20 张照片，系统训练个人模型，再通过固定模板和固定生成链路输出写真。它把 AI 能力用得很好，但仍然通过限制用户自由度来换稳定效果。
\end{knowledgebox}

\subsection{真、像、美}

前面说妙鸭链路固定，但固定链路不等于简单。它真正的难点，是把一个很窄的场景做到足够强的效果函数。

妙鸭的效果标准是三件事。真，是看不出 AI 生成；像，是足够像用户本人，让传播与“我”有关；美，是比真实自己更好看一点。这个三角解释了为什么它能成为爆款：它不是展示 AI 技术，而是把 AI 技术变成一个强烈的消费体验。

\begin{importantbox}{妙鸭的产品教训}
AI 爆款不一定来自最通用的模型，而可能来自一个垂直场景中的明确效果函数。妙鸭的效果函数就是“真、像、美”。
\end{importantbox}

\subsection{为什么这仍是互联网产品}

效果函数解释了妙鸭为什么成功，本节则解释它为什么仍不是 AI Native。关键不在结果是否惊艳，而在产品是否仍然依赖预设流程。

妙鸭本质上仍是流程设计：用户按预设步骤上传照片，选择模板，系统沿固定工艺链生成结果。用户输入不开放，输出也不开放。产品经理能够设计完整流程，工程师按流程实现，体验可控。这是互联网产品方法论的胜利，而不是 AI Native 方法论的完全展开。

\subsection{本章小结}

妙鸭说明，AI 能力用得好，不等于产品范式被 AI 改写。它是一个极好的“AI 加互联网产品”，但不是张月光后来理解的 AI Native。

